\documentclass[10pt,a4paper]{amsart}
\usepackage[margin=19mm]{geometry}
\usepackage{amsmath,amssymb,mathtools}
\usepackage[scr=rsfs,cal=euler]{mathalfa}
\usepackage{xfrac}
\usepackage[unicode,hidelinks,bookmarks=false]{hyperref}
\newcommand{\ie}{i.e.\ }
\newcommand{\cf}{cf.\ }
\newcommand{\resp}{respectively\ }
\newcommand{\Subsection}{\subsection}
\providecommand{\nobreakdash}{\nobreak\hbox{-}}
\setlength{\emergencystretch}{2em}
\multlinegap0pt
\usepackage{tikz}
\usetikzlibrary{cd}
\usepackage{amssymb}
\usepackage{tikz}
\newtheorem{prop}{Proposition}
\newcommand{\Ker}{\operatorname{Ker}\nolimits}
\title{Some Remarks about Integral Categories}
\author{Yaroslav Kopylov, Max Zinchenko}
\date{Source: arXiv:2606.21248v1 (CC BY 4.0)}
\begin{document}
\maketitle

\noindent\emph{Source and licence.} Some Remarks about Integral Categories. Version arXiv:2606.21248v1 (CC BY 4.0), distributed by arXiv under the Creative Commons Attribution 4.0 International licence, http://creativecommons.org/licenses/by/4.0/, as recorded in the arXiv licence metadata for this exact version and shown on the abstract page https://arxiv.org/abs/2606.21248v1. Full licence notice: \texttt{licenses/arxiv-2606.21248-CC-BY-4.0.txt}.\par\smallskip

\noindent\textit{Editorial scope.} Counted unit: the article's prop sa-int2 (source0.tex lines 208--212). The article's complete original proof of it is delivered with it (source0.tex lines 214--246, one \texttt{proof} environment of 33 lines). No mathematics is written, repaired or re-derived: the statement and the proof are the article's own lines, in the article's own notation, and no retained line is substituted or re-flowed.  No further source-local context is needed: every cross-reference of every retained line resolves inside the two delivered spans.  Nothing was omitted from any retained span, so the package carries no omission record.  The retained lines cite 1 works, whose entries are transcribed verbatim from the article's own bibliography source in the e-print and delivered with short editorial labels recorded in the item's reference mapping.  The retained lines contain the article's own diagram code, which the wrapper typesets with the standard tikz packages; no image file is delivered and the package declares no asset dependency.  Editorial formatting only, in addition: an AMS \texttt{amsart} preamble that restates the article's own declarations and macros used by the retained lines, namely the environments \texttt{prop} on separate counters, together with the standard packages those lines need.  The retained environments print in this document as Proposition 1 in order of appearance, whereas the article prints the same items with the numbering of its own full document.  The complete native source of the article as distributed in the arXiv e-print for this exact version is bundled unchanged as \texttt{sources/203260/source0.tex}.
	\begin{prop}\label{sa-int2}
		A~right semi-abelian category $\mathcal{A}$ is right integral if and only if pushouts of monomorphisms along epimorphisms in~$\mathcal{A}$ are monomorphisms.
		
		Dually, a~left semi-abelian category $\mathcal{A}$ is left integral if and only if pullbacks of epimorphisms along monomorphisms in~$\mathcal{A}$ are epimorphisms.
	\end{prop}

	\begin{proof}
		We prove the~first assertion.
		
		If $\mathcal{A}$ is right integral then the~assertion is obvious.
		
		Suppose that pushouts of monomorphisms along epimorphisms in~$\mathcal{A}$ are monomorphisms and consider a~pushout
		$$
		\begin{tikzcd}
			C \arrow[r, "g"] \arrow[d, "\alpha"'] & D \arrow[d, "\beta"] \\
			A \arrow[r, "f"'] & B
			\arrow[from=1-1, to=2-2, phantom, "\scriptstyle\text{PO}" description]
		\end{tikzcd}
		$$
		with $g$ a~monomorphism. Since $\mathcal{A}$ is right semi-abelian,
		$\overline{\alpha}$ is an~epimorphism, and so $\alpha=\alpha_1 \alpha_0$, where $\alpha_1:A'\to A$ is a~kernel and $\alpha_0:C\to A'$ is an~epimorphism. Consider the~pushout
		$$
		\begin{tikzcd}
			C \arrow[r, "g"] \arrow[d, "\alpha_0"'] & D \arrow[d, "\beta_0"] \\
			A' \arrow[r, "f_0"'] & B'
			\arrow[from=1-1, to=2-2, phantom, "\scriptstyle\text{PO}" description]
		\end{tikzcd}
		$$
		Since $f\alpha_1\alpha_0 = f\alpha = \beta g$ and the~above square
		is a~pushout, there exists a~unique morphism $\beta_1:B'\to B$ such that $f\alpha_1 = \beta_1 f_0$ and $\beta=\beta_1\beta_0$. Since $\alpha_0$ is an~epimorphism, it is known and easy to check that
		$$
		\begin{tikzcd}
			A' \arrow[r, "f_0"] \arrow[d, "\alpha_1"'] & A \arrow[d, "\beta_1"] \\
			A \arrow[r, "f"'] & B
			\arrow[from=1-1, to=2-2, phantom, "\scriptstyle\text{PO}" description]
		\end{tikzcd}
		$$
		is also a~pushout. Since $\mathcal{A}$ is right semi-abelian and $\alpha_1$ is a~kernel, we conclude that the~last square is a~pullback (see \cite[Proposition~3.1]{KW}). Thus, $0=\Ker f_0\equiv \Ker f$, and $f$ is a~monomorphism. 
	\end{proof}

\begin{thebibliography}{99}
\bibitem[KW]{KW} Ya. Kopylov and S.-A. Wegner, \emph{On the notion of a semi-abelian category in the sense of {P}alamodov}, Appl. Categor. Struct. \textbf{20} (2012), no.~5, 531--541.
\end{thebibliography}
\end{document}
